\section{Introduction}
\label{sec:intro}

Large language models (LLMs) do not always tell the truth \cite{park2023deception}. Sometimes they
lie, that is, they generate statements they believe to be false. Following Liars' Bench
\cite{kretschmar2025liars}, we adopt an assertion based definition of lying: the assistant lies if it
states something it believes to be false. Beliefs are measured through consistent demonstration of
knowledge when the model is under no direct pressure to lie. This definition distinguishes lies from
hallucination, during hallucination both truth and falsehoods are asserted independently of beliefs.
Prior work found models to be lying in various scenarios, including lying in order to conceal their
prior misconduct \cite{scheurer2023strategic,jarviniemi2024uncovering}, and lying in order to fake
their alignment to seem cooperative \cite{greenblatt2024alignment}. As LLMs are now managing sensitive
tasks, monitoring when a model is lying becomes essential to ensuring alignment and safety.

Two broad families of lie detectors exist: Black-box and White-box. Black-box detectors can only see
the chat transcript, a LLM judge is instructed to read the transcript and evaluate if the model lied
or not \cite{pacchiardi2023catch}. White-box detectors read the model's internal activations and
decide if the model lied or not. The simplest white-box detector is a linear probe, which is a linear
classifier trained on the activations collected from the residual stream
\cite{alain2016understanding,azaria2023internal,marks2023geometry}. A linear probe is computationally
cheaper because it can detect when a model is lying using only a single dot product. Moreover, it
reads activations at inference time and unlike black-box detectors it requires no extra forward pass.
As it reads models internal representations, it can still detect lies even if there is no direct
trace of lying in the transcript \cite{goldowskydill2025detecting}.

The weakness of the probe is generalisation. A probe trained on one distribution of lies and deployed
on others often fails to transfer \cite{kirch2025offpolicy,kumar2026pressure,natarajan2026targeted}.
Liars' Bench makes the failure concrete \cite{kretschmar2025liars}. It contains \V{lb-examples} lies
and honest responses, across seven datasets it captures distinct types of lies generated on-policy by
four popular open-weight LLMs. Those datasets vary on two axes: object of belief and reason for lying.
Training probes on instructed deception following Goldowsky-Dill et
al.~\cite{goldowskydill2025detecting}, reaches balanced accuracy near chance on most of them. The
hardest is the harm pressure knowledge report (\hpkr). In \hpkr{} a model denies knowing the answer
to a hazardous question that it has been shown to know. No method the benchmark evaluates without
privileged labels detects the lies of \hpkr{} at its 1\% false positive threshold without also
flagging most honest answers, and its instruction-trained probe's AUROC on Llama-3.3-70B-Instruct is
\V{lb-mean-auroc-hpkr}, which is worse than chance. That inversion is measured at layer~20, the layer
the benchmark reads; at layer~40, the layer we report, a probe trained the same way ranks \hpkr{} well
but still flags none of its lies (\cref{sec:problem-measured}). Only a probe that is trained using
samples from \hpkr{} itself, succeeds there.

Mahadik and Skapars~\cite{mahadik2026persona} propose using persona coordinates for better
generalisation of probes. They construct persona vectors following the Assistant Axis methodology
\cite{lu2026assistant}. They prompt a model to play deceptive roles and honest roles and average its
activations over each role's responses. It gives one persona vector per role. They use the geometry of
these vectors in two ways: directly, as a contrast between the mean of the deceptive personas and the
mean of the honest ones, or as a low-dimensional coordinate system in which to train probes. Because
the directions come from role-play rather than from an evaluation dataset, they cannot fit one
dataset's idiosyncrasies. On Llama~3.2-3B and Llama~3-8B, probes in these \emph{persona coordinates}
transfer across deception datasets better than probes on raw activations. The authors name larger
models and richer persona sets as future work. We take up both.

When the published persona contrast is rebuilt for Llama-3.3-70B-Instruct and applied to Liars'
Bench, it fails on \hpkr{} in the worst possible way: it \emph{inverts}, ranking the model's lies as
more honest than its honest answers, with AUROC \V{zs-dec-hpkr-L20} at layer~20. In \hpkr{} every lie
is a \emph{denial of knowledge}, while every published deceptive persona lies by \emph{controlling
disclosure}: timing, emphasis and omission. Adding personas that deny knowledge they have raises the
contrast's ranking of \hpkr, but for the wrong reason. The new direction flags nearly every response
that says ``I do not know'', true or false: at layer~40 it flags \V{th-k-deight-flag-hden-L40}\% of
the honest denials. Where the lies and the honest answers are the same sentence, every persona
direction is within reach of random directions, and on Gender-Secret, a model fine-tuned to deny
knowing its user's gender, every persona direction is at chance (\cref{sec:hpkr-tests}). A persona
direction reads the register of a lie, not what makes it a lie.

The definition of a lie offers a second route. A lie is a statement the speaker believes to be false,
so a lie about knowledge is a denial by a model that knows. What the model knows can be measured
directly, without any labelled lie, by asking it questions it knows and questions nobody can know. We
therefore ask whether a lie detector for a 70B-parameter language model, built without any labelled
lies, can separate lies from honest answers when the two use the same words, as lies about knowledge
often do, while keeping the cross-dataset transfer of persona directions. Four research questions
follow.
\begin{itemize}
\item[] \emph{RQ1: Transfer.} Does a detector built from a persona contrast and a knowledge direction,
  with no labelled lies, transfer across Liars' Bench better than supervised probes?
\item[] \emph{RQ2: The same words.} Does it separate lies from honest answers that use the same words,
  as lies about knowledge often do, and does it detect the lies of a model fine-tuned to hide what it
  knows?
\item[] \emph{RQ3: The two parts.} What does each part read: does the persona contrast respond to how a
  response is framed, and does the knowledge direction read whether the model knows the answer while
  the words of its reply stay fixed?
\item[] \emph{RQ4: Robustness and operating points.} Does the result survive resampling the questions,
  new questions and other layers, and how many lies does the detector catch at strict thresholds?
\end{itemize}

In this work, we make the following contributions:
\begin{itemize}
\item \emph{A label-free lie detector.} The extended persona contrast plus the knowledge direction
  transfers better than supervised probes, and separates lies that use the same words as honest
  answers (\cref{sec:transfer,sec:selection}).
\item \emph{A knowledge direction.} It reads what the model knows while the words stay fixed. It is
  built from questions the model knows and questions nobody can know, without labelled lies
  (\cref{sec:k-design,sec:k-results}).
\item \emph{A persona bank for deception by denial.} We add knowledge-denial roles to the published
  bank and scale the method to a 70B model, where a single persona contrast transfers better than
  supervised probes (\cref{sec:banks,sec:transfer}).
\item \emph{A diagnosis of persona directions.} On lies about knowledge, persona directions read the
  words of a denial rather than the lie, which is why the detector needs the knowledge direction
  (\cref{sec:hpkr-tests}).
\item \emph{Tests on fine-tuned models.} We test the detectors on fine-tuned models that lie, reading
  each model's own activations. The knowledge direction detects Gender-Secret's lies there, while the
  persona directions are near chance (\cref{sec:finetuned-reads,sec:k-results}).
\item \emph{A causal test of the persona part.} Steering along the extended contrast makes the model
  conceal an unfavourable fact far more often, while random directions do not (\cref{sec:steering}).
\end{itemize}
